\documentclass[12pt,a4paper]{article}

\usepackage[T1]{fontenc}
\usepackage[utf8]{inputenc}
\usepackage{lmodern}
\usepackage[polish]{babel}
\usepackage{geometry}
\usepackage{graphicx}
\usepackage{hyperref}
\usepackage{enumitem}
\usepackage{booktabs}
\usepackage{float}
\usepackage{caption}
\usepackage{listings}
\usepackage{xcolor}

\geometry{margin=2.5cm}
\hypersetup{
    colorlinks=true,
    linkcolor=black,
    urlcolor=blue,
    pdftitle={Dokumentacja projektu},
}

\lstset{
    basicstyle=\fontfamily{lmtt}\selectfont\small,
    breaklines=true,
    frame=single,
    backgroundcolor=\color{gray!8}
}

\newcommand{\diagram}[3]{%
    \begin{figure}[H]
        \centering
        \includegraphics[width=0.9\textwidth]{#2}
        \caption{#1}
        \label{fig:#3}
    \end{figure}
}

\title{%
    \textbf{Symulacja pracy Szpitalnego Oddziału Ratunkowego}\\[0.4em]
    \large Dokumentacja projektu\\Programowanie obiektowe
}
\date{\today}

\begin{document}

\maketitle
\tableofcontents
\newpage

\section{Wprowadzenie}

Niniejszy dokument opisuje projekt przygotowany na zaliczenie z przedmiotu
programowanie obiektowe. Głównym założeniem projektu jest zaprojektowanie
i zaimplementowanie aplikacji w języku C++, która w sposób uproszczony
odwzorowuje pracę Szpitalnego Oddziału Ratunkowego (SOR). W implementacji
zastosowano podstawowe mechanizmy programowania obiektowego: dziedziczenie,
polimorfizm, enkapsulację oraz kompozycję obiektów.

Projekt nie stanowi pełnego systemu informatycznego szpitala, lecz model
dydaktyczny. Symulacja generuje losowych pacjentów, przydziela im priorytet
w triażu, kolejkuje ich według pilności i obsługuje w gabinetach lekarskich.
Wyniki zapisywane są do pliku CSV, a następnie można je przeanalizować
w prostym panelu webowym napisanym w Pythonie.

\section{Cel i zakres projektu}

\subsection{Cel główny}

Celem projektu jest stworzenie symulacji dyskretnej, która pozwala obserwować,
jak zmienia się obciążenie oddziału ratunkowego w czasie. Model umożliwia
odpowiedź na kilka prostych pytań badawczych:
\begin{itemize}[noitemsep]
    \item ile pacjentów dociera do SOR-u w zadanym okresie,
    \item jak rośnie kolejka oczekujących,
    \item ile gabinetów jest jednocześnie zajętych,
    \item ile osób zostaje ostatecznie wyleczonych.
\end{itemize}

\subsection{Zakres funkcjonalny}

Program realizuje następujące zadania:
\begin{itemize}[noitemsep]
    \item wczytanie parametrów symulacji z pliku konfiguracyjnego \texttt{config.ini},
    \item losowe generowanie pacjentów o różnych typach klinicznych,
    \item triaż wykonywany przez pielęgniarkę (przypisanie priorytetu),
    \item kolejkowanie pacjentów w poczekalni według priorytetu,
    \item leczenie w gabinetach konsultacyjnych prowadzonych przez lekarzy,
    \item symulację pogorszenia stanu pacjentów oczekujących (zależnie od typu),
    \item logowanie stanu systemu do pliku CSV po każdym kroku czasowym,
    \item wizualizację wyników w panelu Streamlit.
\end{itemize}

\section{Opis problemu i motywacja}

Tematem projektu jest Szpitalny Oddział Ratunkowy, czyli miejsce, w którym
pacjenci przyjmowani są w trybie pilnym. W rzeczywistości personel musi szybko
ocenić stan pacjenta (triaż) i zdecydować, kto wymaga natychmiastowej pomocy.
W modelu projektu proces ten uproszczono do trzech poziomów priorytetu:
\emph{zielony}, \emph{żółty} oraz \emph{czerwony}.

\section{Architektura rozwiązania}

\subsection{Przegląd modułów}

Aplikacja składa się z kilku logicznych części:
\begin{description}
    \item[Symulacja C++] Główna logika programu. Klasa
        \texttt{EmergencyDepartment} koordynuje przebieg symulacji.
    \item[Model danych] Hierarchia klas reprezentujących osoby, pacjentów,
        pracowników medycznych oraz pomieszczenia.
    \item[Konfiguracja] Plik INI z czasem symulacji, krokiem czasowym, liczbą
        gabinetów i prawdopodobieństwem pojawienia się pacjenta.
    \item[Logger CSV] Zapis metryk: minuta, zajęte gabinety, liczba oczekujących,
        łączna liczba wyleczonych.
    \item[Panel analityczny (Python)] Skrypt \texttt{dashboard.py} wczytuje CSV
        i rysuje wykresy w Streamlit.
    \item[Testy jednostkowe] Testy Google Test sprawdzają wybrane klasy.
    \item[Dokumentacja API] Dokumentacja kodu generowana przez Doxygen.
\end{description}

\subsection{Przebieg symulacji}

Symulacja działa w pętli czasowej. W każdym kroku (domyślnie co kilka minut
wirtualnych):
\begin{enumerate}[noitemsep]
    \item system próbuje wygenerować nowego pacjenta (z zadanym prawdopodobieństwem),
    \item pielęgniarka wykonuje triaż i pacjent trafia do poczekalni,
    \item aktualizowany jest stan pacjentów oczekujących,
    \item w gabinetach kontynuowane jest leczenie,
    \item wolne gabinety z przypisanym lekarzem przyjmują kolejnych pacjentów
        z poczekalni,
    \item logger zapisuje bieżące statystyki.
\end{enumerate}

Kolejka w poczekalni jest implementowana jako kopiec priorytetowy. Pacjenci
o wyższym priorytecie są obsługiwani wcześniej.

\subsection{Hierarchia klas}

Poniżej przedstawiono najważniejsze relacje między klasami (szczegóły
w diagramach w dalszej części dokumentu):

\begin{itemize}[noitemsep]
    \item \texttt{Person}: klasa bazowa dla ludzi w systemie.
    \item \texttt{Patient}: pacjent; klasa abstrakcyjna z metodami wirtualnymi
        \texttt{performWaitingStep()} oraz \texttt{getDiagnosis()}.
    \item \texttt{StandardPatient}, \texttt{CardiacPatient}, \texttt{TraumaPatient}:
        konkretne typy pacjentów.
    \item \texttt{Employee}: pracownik; klasa abstrakcyjna z metodą
        \texttt{isBusy()}.
    \item \texttt{Doctor}, \texttt{Nurse}: role medyczne.
    \item \texttt{Room}: pomieszczenie; bazowa dla \texttt{WaitingRoom}
        i \texttt{ConsultationRoom}.
    \item \texttt{EmergencyDepartment}: agreguje poczekalnię, gabinety,
        personel i logger.
    \item \texttt{ConfigParser}, \texttt{LoggerCSV}: klasy pomocnicze.
\end{itemize}

\section{Diagramy}

\diagram{Diagram klas systemu symulacji SOR}{img/class.png}{diagram-klas}

\diagram{Diagram obiektów (przykładowa konfiguracja w czasie triażu)}{img/object.png}{diagram-obiektow}

\diagram{Diagram sekwencji: przyjęcie i obsługa pacjenta}{img/sequence.png}{diagram-sekwencji}


\section{Zastosowane mechanizmy obiektowe}

W projekcie wykorzystano kilka koncepcji omawianych na zajęciach:

\begin{itemize}
    \item \textbf{Dziedziczenie}: na przykład \texttt{Patient} dziedziczy po
        \texttt{Person}, a \texttt{CardiacPatient} po \texttt{Patient}.
    \item \textbf{Polimorfizm }: różne typy pacjentów inaczej
        zachowują się podczas oczekiwania (\texttt{performWaitingStep}).
    \item \textbf{Abstrakcja}: klasy bazowe definiują interfejs, a szczegóły
        implementują klasy pochodne.
    \item \textbf{Kompozycja i agregacja}: \texttt{EmergencyDepartment}
        zarządza obiektami poczekalni, gabinetów i pracowników.
    \item \textbf{Enkapsulacja}: pola klas są ukryte w sekcji
        \texttt{private}/\texttt{protected}, dostęp zapewniają metody.
    \item \textbf{RTTI}: w triażu pielęgniarka wyszukiwana jest przez
        \texttt{dynamic\_cast\textless Nurse*\textgreater}.
\end{itemize}

\section{Konfiguracja symulacji}

Parametry wczytywane są z pliku \texttt{config/config.ini}. Przykładowa zawartość:

\begin{lstlisting}
total_time=2139
step_duration=5
spawn_chance=67
num_rooms=2
\end{lstlisting}

\begin{table}[H]
    \centering
    \caption{Opis parametrów konfiguracyjnych}
    \begin{tabular}{@{}llp{7cm}@{}}
        \toprule
        Parametr & Typ & Znaczenie \\
        \midrule
        \texttt{total\_time} & int & Czas trwania symulacji w minutach wirtualnych \\
        \texttt{step\_duration} & int & Długość jednego kroku symulacji \\
        \texttt{spawn\_chance} & int & Szansa (od 0 do 99) pojawienia się pacjenta w kroku \\
        \texttt{num\_rooms} & int & Liczba gabinetów konsultacyjnych \\
        \bottomrule
    \end{tabular}
\end{table}

\section{Wyniki symulacji i panel analityczny}

Po uruchomieniu program zapisuje dane do \texttt{output/simulation.csv}. Każdy wiersz
zawiera m.in.\ bieżącą minutę, liczbę zajętych gabinetów, pacjentów oczekujących
oraz skumulowaną liczbę wyleczonych osób.

Panel \texttt{dashboard.py} (Streamlit + pandas) umożliwia szybki podgląd wykresów
bez ręcznej analizy pliku. Pozwala to ocenić, czy na przykład zwiększenie liczby
gabinetów skraca kolejkę po zmianie konfiguracji i ponownym uruchomieniu symulacji.

\section{Wykorzystane narzędzia i technologie}

\begin{table}[H]
    \centering
    \caption{Stos technologiczny projektu}
    \begin{tabular}{@{}lp{9cm}@{}}
        \toprule
        Narzędzie & Zastosowanie \\
        \midrule
        C++ & Implementacja silnika symulacji \\
        GCC / g++ & Kompilator \\
        GNU Make & Automatyzacja budowania \\
        Google Test & Testy jednostkowe \\
        Python 3 & Panel wizualizacji wyników \\
        Streamlit, pandas & Interfejs webowy i analiza niezmiernie użytecznych danych CSV \\
        Doxygen & Automatyczna dokumentacja kodu źródłowego \\
        \LaTeX & Niniejsza dokumentacja projektowa \\
        Nix  & Reprodukowalne środowisko budowania i paczki \\
        \bottomrule
    \end{tabular}
\end{table}

\section{Budowanie i uruchamianie}

\subsection{Kompilacja i uruchomienie symulacji}

W katalogu projektu, po wejściu do powłoki Nix (\texttt{nix develop}), symulację
uruchamia się poleceniami:

\begin{lstlisting}
make
make run
\end{lstlisting}

Plik wykonywalny powstaje w \texttt{bin/symulacja}. Przed uruchomieniem warto
sprawdzić parametry w \texttt{config/config.ini}.

\subsection{Testy}

Testy uruchamiane są poleceniem:

\begin{lstlisting}
make test
\end{lstlisting}

\subsection{Panel analityczny}

Po wygenerowaniu pliku CSV:

\begin{lstlisting}
make dashboard
\end{lstlisting}

\subsection{Dokumentacja}

Cała dokumentacja (Doxygen + PDF) budowana jest jedną paczką Nix:

\begin{lstlisting}
nix build .#docs
\end{lstlisting}

Wynik instalacji zawiera:
\begin{itemize}[noitemsep]
    \item \texttt{result/doxygen/}: dokumentacja HTML z Doxygen,
    \item \texttt{result/docs.pdf}: niniejszy dokument w formacie PDF.
\end{itemize}

\end{document}